\bookchapter{Bab 47: Buku Rujukan Pantas dan Sumber}

\chapterintro{Bab ini ialah buku rujukan pantas yang boleh digunakan secara berasingan. Apabila anda tidak ingat sesuatu prinsip atau perlu membuat keputusan pantas, teruskan membaca dari sini. Digalakkan menjadikannya “kad poket”: rujukan pantas lima prinsip utama, pokok keputusan alat, soalan lazim dan sumber pembelajaran — tiada satu halaman pun tertinggal.}

\booksection{Kad rujukan pantas prinsip pemikiran}

\begin{center}
\begin{tabularx}{\linewidth}{|X|X|X|}
\hline
\textbf{Prinsip} & \textbf{Dalam satu ayat} & \textbf{Perkara utama pelaksanaan} \tabularnewline
\hline
Think Before Coding & Fikir dahulu dengan jelas sebelum bercakap & Jawab tiga soalan: buat apa? Bagaimana rupa hasilnya? Ada kekangan apa? \tabularnewline
Simplicity First & Selesaikan masalah dengan kod paling sedikit & Setiap kali hanya buat satu fungsi; buang kerumitan yang tidak perlu \tabularnewline
Surgical Changes & Ubah suai dengan tepat tanpa melampaui batas & Nyatakan skop ubah suai dengan jelas; tolak "sambil tukar sedikit" daripada AI \tabularnewline
Goal-Driven Execution & Sebut "buat apa" bukan "cara buat" & Beri sasaran, bukan langkah; biarkan AI memilih laluannya sendiri \tabularnewline
Iterate \& Verify & Kitaran jana-sahkan-baiki & Pengesahan tiga lapisan: fungsi, sempadan, seni bina \tabularnewline
\hline
\end{tabularx}
\end{center}

\begin{center}\small\color{gray}Kad rujukan pantas lima prinsip pemikiran utama\end{center}

\booksection{Pokok keputusan pemilihan alat}

Apabila anda perlu memilih alat pengaturcaraan AI, pilih mengikut laluan keputusan berikut. Perhatian: ini bukan senarai kedudukan “mana yang terbaik”, tetapi rangka kerja penilaian “mana yang paling sesuai untuk anda”.

Anda mahu buat apa?

Sifar asas, mahu rasa cepat pilih Builder awan (v0/Bolt/Replit)

Pembangunan harian, perlukan pelengkap kod pilih AI terbina dalam IDE (Cursor/Trae)

Projek besar, perlu operasi pukal pilih Agent terminal (Claude Code)

Uji idea cepat dengan prototaip pilih Builder awan atau Trae SOLO

\begin{center}\small\color{gray}Rajah 6: Pokok keputusan pemilihan alat pengaturcaraan AI\end{center}

\booksection{Soalan lazim (FAQ)}

Soalan 1: Saya langsung tidak tahu pengaturcaraan, bolehkah saya menggunakan pengaturcaraan AI?

Boleh. Pada 2026, alat pengaturcaraan AI sudah cukup kuat; anda boleh sepenuhnya menerangkan keperluan dalam bahasa semula jadi dan membiarkan AI menjana kod. Tetapi ada dua cadangan: pertama, mulakan dengan alat Builder awan (seperti v0 dan Bolt) yang tidak memerlukan persediaan persekitaran pembangunan; kedua, pelajari konsep asas sambil menggunakan — anda tidak perlu mahir sintaks, tetapi anda perlu dapat “memahami” maksud umum kod yang ditulis AI, barulah penerimaan dapat dilakukan dengan berkesan.

Soalan 2: Bahasa pengaturcaraan yang mana patut saya pelajari?

Bergantung pada apa yang anda mahu buat. Frontend laman web — JavaScript/TypeScript; pemprosesan data atau AI — Python; backend peringkat perusahaan — Java atau Go. Tetapi lebih penting: pada era pengaturcaraan AI, kepentingan “bahasa” sedang menurun. AI boleh membantu anda mengendalikan butiran sintaks; yang perlu anda fahami ialah “struktur program” — apa itu fungsi, apa itu gelung, apa itu struktur data. Konsep-konsep ini melangkaui bahasa; sekali difahami, bertukar ke bahasa mana pun menjadi cepat.

Soalan 3: Adakah kod yang ditulis AI ada isu hak cipta?

Sehingga Ogos 2026, ini masih satu isu undang-undang yang dipertikaikan. Secara umumnya, kod yang dijana AI tidak dianggap mempunyai pemilik hak cipta yang jelas dalam kebanyakan bidang kuasa. Tetapi jika anda menggunakannya untuk projek komersial, cadangannya: (1) ketahui klausa harta intelek dalam terma perkhidmatan alat AI yang anda gunakan; (2) buat semakan dan ubah suai manusia ke atas kod penting supaya mengandungi sumbangan kreatif anda; (3) elakkan membiarkan AI menyalin terus kod sumber terbuka yang mungkin ada hak cipta.

Soalan 4: Adakah AI akan menggantikan pengaturcara?

Dalam jangka pendek, tidak. Yang digantikan oleh AI bukan kerjaya “pengaturcara”, tetapi kerja “mengekod baris demi baris”. Sama seperti kalkulator tidak menggantikan ahli matematik, malah membebaskan ahli matematik daripada pengiraan tangan yang menjemukan. Alat pengaturcaraan AI mengalihkan pengaturcara daripada “menulis kod” kepada “memikirkan masalah” — ini sebenarnya meningkatkan had nilai pengaturcara. Daya saing masa depan bukan pada “boleh menulis berapa banyak baris kod”, tetapi pada “boleh menyelesaikan masalah seberapa kompleks”.

Soalan 5: Bagaimana menilai sama ada output AI itu betul?

Kaedah pengesahan tiga lapisan. Lapisan pertama: pengesahan fungsi — bolehkah ia berjalan? Betulkah input dan outputnya? Ini yang paling asas. Lapisan kedua: pengesahan sempadan — apa yang berlaku dengan input kosong, input terlalu panjang atau input luar biasa? Ini dapat menemui masalah tersembunyi. Lapisan ketiga: pengesahan seni bina — adakah struktur kod munasabah? Adakah ada yang dikodkan terus? Adakah mudah dikembangkan? Ini memerlukan pertimbangan teknikal tertentu daripada anda. Jika anda tidak dapat menilai lapisan ketiga, sekurang-kurangnya lakukan dua lapisan pertama — boleh berjalan dan sempadan normal sudah mengatasi kebanyakan pemula.

\booksection{Sumber pembelajaran yang disyorkan}

Berikut ialah sumber pembelajaran yang disyorkan sehingga Ogos 2026. Sila ambil perhatian, pautan sumber internet boleh berubah; jika sesuatu pautan tidak lagi berfungsi, mencari tajuknya biasanya dapat menemui versi terkini.

\begin{center}
\begin{tabularx}{\linewidth}{|X|X|X|}
\hline
\textbf{Jenis} & \textbf{Sumber} & \textbf{Sebab disyorkan} \tabularnewline
\hline
Bahan asal & CLAUDE.md milik Andrej Karpathy & Punca lima prinsip utama; hanya 70 baris tetapi setiap patah kata bernilai \tabularnewline
Bahan asal & Dokumentasi rasmi Anthropic (Claude Code) & Panduan amalan terbaik bagi Agent terminal \tabularnewline
Tutorial & Koleksi Skills yang disusun komuniti Vibe Coding & Kumpulan teknik praktikal yang dirumuskan oleh pengamal komuniti \tabularnewline
Tutorial & Dokumentasi rasmi Cursor dan saluran YouTube & Panduan berautoriti menggunakan AI terbina dalam IDE \tabularnewline
Komuniti & awesome-claude-code di GitHub & Amalan terbaik dan koleksi plugin Claude Code \tabularnewline
Komuniti & Komuniti r/cursor di Reddit & Perkongsian pengalaman dan perbincangan pengguna Cursor seluruh dunia \tabularnewline
Asas & MDN Web Docs (developer.mozilla.org) & Dokumentasi rujukan paling berautoriti untuk pembangunan frontend \tabularnewline
Asas & Tutorial rasmi Python (docs.python.org) & Titik mula rasmi untuk belajar Python \tabularnewline
\hline
\end{tabularx}
\end{center}

\begin{center}\small\color{gray}Senarai sumber pembelajaran yang disyorkan\end{center}

\booksection{Penutup: alat di tangan, metodologi di hati}

Kembali kepada kepercayaan teras buku ini: alat beriterasi setiap tiga bulan, metodologi hanya berevolusi setiap tiga tahun. Pada hari ini, Ogos 2026, keupayaan alat pengaturcaraan AI sudah kuat sehingga menakjubkan — anda boleh mencipta satu laman web dengan satu ayat, menjana satu aplikasi full-stack dengan satu perenggan penerangan, dan menyiapkan satu refaktor kompleks dengan satu arahan. Tetapi semakin kuat alat, semakin terserlah pentingnya metodologi — kerana kekuatan alat membesarkan setiap keputusan anda.

Jika keperluan anda jelas, AI dapat membantu anda melaksanakannya dengan cekap; jika keperluan anda kabur, AI juga cekap menghasilkan benda yang salah. Jika anda mematuhi kesederhanaan diutamakan, AI dapat membantu anda mengesahkan idea dengan pantas; jika anda tamak untuk menyeluruh, AI juga pantas menimbun hutang teknikal. Jika anda mengubah suai dengan tepat, AI dapat mengoptimumkan kod secara presisi seperti pembedahan; jika arahan anda kabur, AI juga pantas memperkenalkan bug secara besar-besaran.

Inilah sebab buku ini menggunakan sebahagian besar ruangannya untuk membincangkan metodologi dan bukan operasi alat. Bukan kerana alat tidak penting, tetapi kerana metodologi ialah “pendarab” alat — semakin baik metodologi, semakin besar nilai alat; jika metodologi sifar, alat sekukur mana pun hanyalah sifar didarab dengan apa jua nombor.

Anda tidak lagi mengekod bersendirian. Anda seorang panglima, AI ialah tentera anda. Kualiti pemikiran anda menentukan kekuatan tempur tentera itu. Mulai hari ini, persenjatai pemikiran dengan metodologi dan perbesar keupayaan dengan alat—itulah keseluruhan rahasia pengekodan AI.

\booksection{Rujukan}

\begin{enumerate}
  \item Andrej Karpathy. (2025). Vibe Coding. Social media post introducing the concept of vibe-driven programming with AI.
  \item Andrej Karpathy. (2025). CLAUDE.md - Project guidelines for AI-assisted development. GitHub repository, 70+K stars.
  \item Anthropic. (2026). Claude Code: Terminal-native AI programming. Official documentation and best practices.
  \item Cursor Team. (2026). Cursor 3: The next generation of AI-first IDE. Release notes and documentation.
  \item OpenAI. (2026). Codex CLI: Bringing AI to the terminal. Technical documentation.
  \item Model Context Protocol (MCP) Specification. (2025-2026). Open protocol for AI-tool integration.
  \item Vibe Coding Community. (2026). Community-curated skills and best practices for AI-assisted programming.
  \item Spec-Driven Development (SDD) methodology. (2026). Emerging practices in specification-first AI programming.
\end{enumerate}
